\chapter{The task and its environments}
\label{ch:task}

\section{The canonical cued change-detection task}
The base environment is a spatially cued orientation-change-detection task built to
reproduce the essential structure of primate covert-attention experiments. Each trial is
a short video of $50\times50$ grayscale frames. The canonical (paper) timeline is
\textbf{seven frames}:
\begin{itemize}
\item $t=0$: black screen.
\item $t=1$: a spatial \emph{cue} appears at one stimulus location --- top-left ($S_1$) or
bottom-right ($S_4$). The cue is a disc whose surrounding ring is completed in proportion
to the cue's \emph{validity} (the probability that, on a change trial, the change will
occur at the cued location): $25\%$, $50\%$, $75\%$, or $100\%$.
\item $t=2$: black screen (the cue is removed; its information must now be held in memory).
\item $t=3,4$: four Gabor-like stimuli appear, one per quadrant ($S_1$ top-left, $S_2$
bottom-left, $S_3$ top-right, $S_4$ bottom-right), at random orientations at least
$45^\circ$ apart, with per-frame orientation ``wobble''.
\item $t=5$: on half of trials (\emph{change trials}) one stimulus's orientation steps by
$\Delta$ degrees; on the other half nothing changes. The change persists to $t=6$.
\item $t=6$: final frame.
\end{itemize}
At every frame the agent chooses \emph{wait} or \emph{declare change}. Declaring ends the
trial. A correct declaration at $t\geq5$ earns reward $1$; waiting out a no-change trial to
$t=6$ also earns $1$; everything else earns $0$. The two errors are therefore misses
(failing to declare a real change) and false alarms (declaring on a no-change trial or
declaring prematurely).

\section{Difficulty control}
Difficulty is set by two knobs. \textbf{Orientation noise} $\sigma$ adds frame-to-frame
jitter $\delta_{it}\sim\mathcal N(0,\sigma)$ to each stimulus orientation; $\sigma=5$ in
the canonical task. The \textbf{change magnitude} $\Delta$ is drawn per change trial from
$\mathcal U(-k,k)$, with $k$ starting at $65^\circ$ and shrinking by a curriculum (factor
$0.9$ each time rolling accuracy exceeds $0.85$, floor $8^\circ$) so the task hardens as
the agent improves. For psychometric analysis we override $\Delta$ to fixed magnitudes and
freeze the curriculum, exactly as one would titrate difficulty in an animal.

\section{Why the seven-step task is hard for reinforcement learning}
The short horizon makes the task deceptively difficult to \emph{learn} even though it is
easy to \emph{solve} (a scripted ``wait until the change is visible, then press'' policy
scores $1.000$). Two degenerate strategies sit at $0.5$ accuracy and act as attractors:
\emph{always-wait} (correct on every no-change trial, misses every change) and
\emph{always-press-at-the-deadline} (catches every change, false-alarms every no-change).
Because the change persists for only two frames, the reward-bearing ``correct press''
window is narrow, so the discriminating policy is hard to discover by exploration, and the
policy entropy collapses before discovery. This is the central learning obstacle the
perception and architecture findings of later chapters are really about.

\section{The task battery}
The same environment generalises to the standard NHP battery by editing the reward branch
and the cue, which is how James proposes to broaden the empirical paper's scope. The
implemented and planned variants are:
\begin{itemize}
\item \textbf{\code{validity4}} --- the canonical four-stimulus, uniform-reward,
validity-driven task above. The workhorse of this report.
\item \textbf{\code{vda4}} --- value-cue variant: the cue colour (red/green/blue) sets the
reward magnitude ($5/3/1$), so behaviour can be read for value direction as well as
spatial cueing.
\item \textbf{\code{setsize\{2,4,9\}}} --- set-size manipulation; the nine-stimulus
($3\times3$) version with five validity levels (including an uninformative $0$) tests
whether the cueing benefit \emph{grows} with set size, the bridge to the normative
set-size/accidental-hit theory (Paper B).
\item \textbf{distractor} (the \code{v7}/\code{v11\_part5} env) --- the cue selects a
\emph{side}; the uncued side independently changes with probability $0.5$ and must be
ignored (pressing on it is a false alarm). Tests genuine selection against a competitor.
\item \textbf{\code{luo\_maunsell}} --- reward-structure manipulation (sensitivity vs.
criterion sessions) after Luo \& Maunsell (2018).
\item \textbf{\code{krauzlis}} --- attend/ignore manipulation after Krauzlis and
colleagues, where an uncued change must be actively ignored.
\end{itemize}
All variants share the clean \code{ChangeDetectionEnv} interface: \code{reset} draws the
cue, validity, and change; \code{step} exposes the two actions and the reward branch. This
uniformity is what makes ``train one architecture across the whole battery'' a tractable
research goal rather than a per-task engineering effort.

\section{The long-horizon variant}
A twenty-nine-step variant (\code{T}=29, change onset drawn uniformly in $[11,25]$) is used
throughout the architecture history. It is the same task with a much longer post-change
window, and as Chapter~\ref{ch:horizon} shows, that single difference --- the length of the
window over which a change can be integrated --- explains most of the otherwise puzzling
discrepancies between the two horizons.
